% part: intuitionistic-logic
% chapter: soundness-completeness
% section: completeness-thm

\documentclass[../../../include/open-logic-section]{subfiles}

\begin{document}

\olfileid[pt]{int}{sc}{cpl}

\olsection{O teorema da completude}

\begin{thm}\ollabel{thm:completeness}
  Se $\Gamma \Entails !A$, então $\Gamma \Proves !A$.
\end{thm}

\begin{proof}
  Provamos a contrapositiva: suponha que $\Gamma \Proves/ !A$. Então, por
  \olref[lin]{lem:lindenbaum}, existe um conjunto primo~$\Gamma^*
  \supseteq \Gamma$ tal que $\Gamma^* \Proves/ !A$. Considere o
  modelo canônico~$\mModel{M(\Gamma^*)}$ para $\Gamma^*$, conforme definido em
  \olref[mod]{defn:canonical-model}. Para todo $!B \in \Gamma$, $\Gamma^*
  \Proves !B$. Observe que $\Gamma^*(\emptyseq) = \Gamma^*$. Pelo Lema da
  Verdade (\olref[tru]{lem:truth}), temos
  $\mSat{M(\Gamma^*)}{!B}[\emptyseq]$ para todo $!B \in \Gamma$ e
  $\mSat/{M(\Gamma^*)}{!A}[\emptyseq]$. Isso mostra que $\Gamma
  \Entails/ !A$.
\end{proof}

\begin{prob}
  Mostre que, se $!A$ contém apenas !!{propositional variable}s, $\lor$
  e~$\land$, então $\Entails/ !A$. Use isso para concluir que $\lif$ não é
  definível na lógica intuicionista a partir de $\lor$ e~$\land$. 
\end{prob}
  
\begin{prob}
  Usando o teorema da completude, prove que, se $\Proves !A \lor !B$, então
  $\Proves !A$ ou $\Proves !B$. (Sugestão: suponha $\mSat/{M_1}{!A}$ e
  $\mSat/{M_2}{!B}$ e construa um novo modelo $\mModel{M}$ tal que
  $\mSat/{M}{!A \lor !B}$.)
\end{prob}
  
\begin{prob} 
  Mostre que, se $\mModel{M}$ é um modelo relacional que usa uma ordem linear, então
  $\mSat{M}{(!A \lif !B)\lor(!B \lif !A)}$.
\end{prob}

\end{document}
